\documentclass[12pt]{article}  
\newcommand{\say}[1]{``#1''}  
\newcommand{\nsay}[1]{`#1'}  
\usepackage{endnotes}  
\newcommand{\B}{\backslash{}}  
\renewcommand{\,}{\textsuperscript{,}}  
\usepackage{setspace}   
\usepackage{tipa}  
\usepackage{hyperref}  
\begin{document}  
\doublespacing  
\section{\href{recovery.html}{On Recovery}}  
First Published: 2026 June 28

\section{Draft 2: 28 June 2026}

Recovery is such a strange concept to me.  
On the one hand, it brings forth for me an image of restoration, returning back to some prior state.  
On the other\footnote{Gardner says I don't need to say hand}, there's no returning to any previous point.  
When something is broken, repair is making something new, especially when it comes to a living being.

Every cell inside of me is being constantly replaced.  
Most every molecule is constantly being reacted and re-made, if not actively filtered out and replaced itself.  
I would have few doubts in my mind that most of the individual atoms themselves\footnote{in as much (inasmuch?) as I still believe that distinct atoms exist within a molecule. Eh. It's a helpful abstraction, at least} are even being cycled through relatively quickly.

So, then, what does recovery mean to a body in constant change?

In this case, my acute recovery is focused on the large patch of skin that was removed from my arm.  
Something that feels somewhat strange now that I've taken a moment to think about it is how our body can regrow each part of our skin at the same place.  
The freckles and moles I have do not change location.  
In some ways this makes sense; the relative location of my skin does not change, so if there's a piece that says to grow weirdly, I suppose I would expect it to keep regrowing weirdly.  
Right now, though, I'm less sure of that fact.

When the skin was removed from my arm, the surgeon\footnote{or someone} sewed the edges together.  
The goal is to minimize the amount of scar tissue that forms.  
That raises an interesting question for me, though: what happens to the relative skin location of each chunk.

There are a few ideas that would make the most sense to me:  
\begin{itemize}

\item The body, upon replacing all the skin cells, reproduces whatever arrangement existed before things are removed.  
Given that the goal of removing skin is often to remove a single set of mutated cells, this may not be what happens.  
Given that I was told to expect my mutation might grow back, though, it's also possible.

\item The body has the mental map of skin cells. There's a distortion, where two previously disparate\footnote{is that the word for disconnected? non-contiguous?} areas are now touching.  
The body is an algorithm, and so it has no relation between radial location of different skin locations.

\item The body has a radial map, and it slowly regenerates the mole from its sewn location back to where it belongs. Removed area is replaced with new skin that maps to the specific radial location.

\end{itemize}

Unfortunately, I did not look for the location of any mole right on the edge of where the incision happened, so I do not know what my body will be treating as the answer to recovery.  
I'm sure that, in general, the body does some combination of the three.  
Honestly, given the strangeness of the human body, I'm sure that it actually does some other thing.

But, that's neither here nor there.  
What is here is the fact that my body is rebuilding itself.  
It's strange to me, therefore, that I am intentionally weakening and restricting my body from maintaining its strength.

It makes sense when I stop to think about how all that will stretch and tear, and I don't want to stretch and tear the stitches or the fragile web of new skin that connects their pieces.  
It's fun to me, though, that to improve one injury, I must allow other parts to become weaker.

This reminds me of what I've heard is a difference between European and Chinese teaching styles.  
Both compare children to fields.  
However, while a field of wheat needs time to rest and years to sit fallow, rice paddies improve when they are constantly maintained.\footnote{is this true? Who can say}  
As a result, European models include long breaks and discrete rest, while Asian models do not.  
I'm sure that this is moreso racial storytelling than anything real, but it's something that can help here; rather than think about this time as becoming weaker, it's time for my body to rest fallow.  
This is the time that I rest to allow for better growth in the future.

Fallow is something that I think we need to get back to as a culture.  
Not necessarily in the initial context of farming, I'm all for man's domination of the natural world, but in the realm of our lived experience.  
Unlike the people who extol boredom and the experience of not having anything to do, I think that it may be better for me and others to consider the idea of letting the mind and body sit fallow.

When a field is left fallow, whatever happens to grow on it grows.  
This allows the minerals and chemicals that desired produce require to refill naturally.  
As we have improved farming techniques, lettings fields fallow has become less important.  
First, we know what plants can be used to replace missing nutrients, and we can intentionally seed them instead.  
Second, we know what chemicals are the nutrients, and we can manually add them to the soil and plants as they need them at each different stage of growth.  
This is certainly better; we have fewer years of famine as a result of fields not producing.

But, there's a value in letting nature take its course.  
I'm sure of that fact, even if the analogy may not work for farming.  
OH wait!

When fields are left to fallow, the plants which return are often plants which local pollinators yearn for.  
Local pollinators, for all that we tend to consider them unneeded, are responsible for large portions of agriculture.  
Giving them opportunities to eat the diets they most desire helps with that.

Second, when fields are scheduled to be fallow, we are reminded that time exists not just within the cycle of the year, but within the cycle of multiple years.  
We have removed the Jubilee year from our society, and we are the weaker for it.  
So much of modern business relies on this idea that there will be eternal debts.

Nothing, however, is eternal.

If every debt came issued with the knowledge that at the next forty-nine or fifty year mark\footnote{although I am not educated, I'm generally now on team forty-nine not fifty year cycle}, what changes would happen?  
Well, it certainly would be harder to get a loan in the years directly before the Jubilee.  
We would not have this issue of generations of wage debt.  
The company store never would have been able to trap people for their entire working life.

I note that, for all that we are consistently being told by the fascists and proto-fascists in power that we need to return to biblical principles, not one has seriously proposed that we have universal debt forgiveness.

When we do not let the ground rest, why would we expect that we would let our fellow man rest?  
The way we treat any part of Creation mirrors the way that we treat all of it.  
In the past, when we needed to let the world rest, we understood that we needed to let one another rest.  
Today, we know what chemicals we need to make someone productive.  
We know exactly when and how to give them for optimal performance.\footnote{or, at least, optimal performance in the metrics we can measure}

There's evidence that fruit trees grow more fruit when they are not monocultured.  
That's more work to harvest, though, so we will not do that.

As I take this time for recovery, I am reminded that the idea that I, a single man, would be expected to keep my home, feed myself, work, and all else is a completely modern idea.  
The boarding house was, so far as I understand it, the most stable iteration of the removal of young men from their families to follow their career.  
And yet, that idea feels so absurd now.

What does it mean to recover?

Recover comes, through steps, from the idea of seizing again.  
It is to recapture something which has been\footnote{presumably} lost.

To recover, then, we must lose something.  
To recover, we must grasp something.

To recover, I must accept that I have lost many of the aspects of me that were my most joyful.

To recover, I must grasp for them.

\section{Draft 1: 28 June 2026}

I've been back home for the past two or so weeks recovering from a surgery that was far more disabling than I expected.  
That was not in my plans for the summer, but I'm also not entirely unhappy with the fact that it happened.  
I love my time with my family, and it's been really nice being able to focus on the recovery and my family life.

It's also been nice to be able to use this time to work on recovering some mental health.  
I think that I've posted here in the past with the fact that my rate of blogging loosely correlates with my mental health, at least over a few week to month term.\footnote{not counting hiatus years, which I think were generally decent}  
Because I have had so few obligations on my time and attention, I have been able to put my focus towards figuring out what serves me best.

Some major takeaways from the physical aspect of recovery:

I hate having a shaved arm. It's very itchy, especially since it's wrapped in such a way that I can't really use my arm.

I miss working out, and especially working out in Pilates classes. They feel really nice in the moment and they help me feel more centered in general. My hips are tight now, which is causing some lower back pain. Probably that's because Pilates helped me strengthen some muscles that were previously weaker, and so now that muscular strength and endurance is degrading\footnote{are degrading?}, there's a mismatch in what my body can do. It will be nice to be able to work out again.

When my arm is splinted, there's still a fair amount of rotation possible in the shoulder. This is bad, as it means that I have to make sure my shoulder is rotated correctly to prevent my arm from hurting.

While I can use my right hand for a task for a few moments, within a minute or so there is a large amount of pain. This reminds me that I do not respect my boundaries or listen to my own body.

I am not great at remembering to take medicine on a routine that's more often than \say{upon waking up}.

Showering is hard when I can't use my right arm. Many things are hard, but the fact that I also have to ensure I keep it dry absolutely doesn't help.

Great!

Some mental notes:

The things I have been given to help with focus help. In particular, when I am on my standard working dose of stimulants\footnote{read: the dose of stimulants I often am taking while at work}, I am able to churn through large amounts of work in relatively short periods of time. When I remove caffeine from the equation, I find that I am less able to go for as long, and I also find that I am less able to stay awake through the day.\footnote{though the fact that we're generally a napping family now also doesn't help me stay up} When I remove all stimulants, I cannot do anything, it feels like.

I often feel like there's this vague cloud of unease around me. When I sit down to journal, the cloud disperses, letting me know that, in retrospect, it was a worry that there might be something on my mind, rather than anything in particular on my mind.

Being in a relationship with my partner is fantastic. It both gives me another reason to feel better and helps me when I'm thinking about the future. I've not made a secret in life or on this blog that one big reason I have yearned for a partner in the past is for a reason to build my life along a certain path.  
In front of me at all times are infinite branching directions I can take my life.  
So many of these branches are now useless to me, as they do not involve my partner or are actively antithetical to remaining with her.  
That's been honestly really nice as I consider my future.

Much as I would like to believe that I am a being that exists separately in separate spheres, I am in fact a composite being.  
Stressors in one part of my life bleed over into others, and when any part of me is struggling, all parts of me struggle more.  
This does also work in reverse, however.  
As any part of me improves, the entirety of me is bettered.

So, that was far more rambly than expected.  
Let's take a bit, call a friend\footnote{did the daily reflection first today}, and then come back and see what we think.  
Great.

Called the friend.\footnote{coming back about fifty minutes later}

\section{Upcoming Posts}  
\begin{itemize}

\item Wonder  
\item Curse of knowledge, esp re. Music Theory\footnote{stolen from a prior list}.

Most of a draft done, though rambly.  
\item Overview of a course on astrochem

I think that I have writing somewhere about it.  
\item RebelFit

A rambly draft done, and some more work which tells me that I was not doing anything wrong in the past, approximation just takes time.  
\item Why aspiring writers should practice typing\footnote{and this should also like tie in the thing that a lot of writers are told to do which is directly write good poetry or prose down to get it in the hand. This also has the effect of practicing hand writing. I think that there's also the whole like \say{when we want to be better at running, we also stretch}}

\end{itemize}

\section{Daily Reflection: 28 June 2026}  
\begin{itemize}  
\item How's things?

Decent! Haven't been keeping up on most things, but that's life sometimes.

\item What's a recent win?

Annotated my entire CV! That's great and exciting.

\item You have the big tapestry embroidery project! How's that going?

Not great. On Friday (two days ago), I did manage to get a few four by four blocks with gaps done, so now I'll go through and fill them in to see what the color looks like.\footnote{Now meaning when I decide to}

\item Do you feel like you're taking good care of yourself?

Eh. I'm being gentler with myself, which is probably good. I was reminded yesterday that caffeine does make a big difference in my productivity, and also that breaks are bad.  
When I take a break, I do manage to get back into the work, but all I want to do next is take another break, and then that second break lasts forever.\footnote{or, at least, it did last time}

\item What could you be doing better?

I had things that I wanted to be doing daily this break, and I haven't been doing all of them as much as I'd like. Could be doing those more.

\item Are you appropriately keeping in touch with those you love?

Mostly! Forgot that I was to call someone yesterday, so will call them after this daily reflection or maybe whole blog post.

\item Big upcoming events? Preparation for them?

\begin{itemize}

\item I have a two week camp equivalent coming up in mid July. I have submitted the necessary paperwork, now I just need to wait to go back to work to remind the work that it's happening in case it didn't get across.  
\item I'm visiting my partner in mid July. I have booked tickets and shared the itinerary.

\end{itemize}

\item Once again, then, how's it going?

Honestly really well. It is really nice to be at home with my family, and I do really appreciate the time that this is giving me to recover and reflect. I do wish that I was being better about journaling on paper, because I do really and legitimately believe that it the thing I do with my life that best serves my mental health in the short and medium term. Despite that fact, though, I am not about to go do so.\footnote{because I am doing this blog post because otherwise I lack faith that it will be done at all}

\end{itemize}

\end{document}